% answers_short.tex —— 2011 年·湖南卷·文科 · 简短答案速查
% 由 tools/make_short_answers.py 自 answers.tex 的 \gksolhead/\gkanswer 提取，
% 并与 2011/2011_数学真题_湖南文_答案要点.md 逐题比对后生成，勿手改（改源 MD 后重跑）。

\begin{gknotice}
\textbf{使用说明}：本节只给最终答案便于快速核对，完整推导见解析版。答案为编写者独立推导，\textbf{非官方标准答案}。
\end{gknotice}

\gkpart{选择题}
\begin{center}\begin{tabular}{ccccccccc}
\toprule
题号 & 1 & 2 & 3 & 4 & 5 & 6 & 7 & 8 \\
\midrule
答案 & B & C & A & C & A & C & B & B \\
\bottomrule
\end{tabular}\end{center}

\gkpart{填空题}
\begin{itemize}[leftmargin=2.6em]
  \item[\textbf{9.}] $2$。
  \item[\textbf{10.}] $40$。
  \item[\textbf{11.}] $\dfrac{15}4$。
  \item[\textbf{12.}] $6$。
  \item[\textbf{13.}] $(-4,-2)$。
  \item[\textbf{14.}] $3$。
  \item[\textbf{15.}] (1) $5$；(2) $\dfrac16$
  \item[\textbf{16.}] (1) $1$；(2) $16$
\end{itemize}

\gkpart{解答题}
\begin{itemize}[leftmargin=2.6em]
  \item[\textbf{17.}] (1) $C=\dfrac\pi4$；(2) 最大值 $2$，此时 $A=\dfrac\pi3$，$B=\dfrac{5\pi}{12}$
  \item[\textbf{18.}] (1) 频率依次为 $\dfrac1{20},\dfrac3{20},\dfrac4{20},\dfrac7{20},\dfrac3{20},\dfrac2{20}$；(2) $\dfrac3{10}$
  \item[\textbf{19.}] (1) 证明见解析；(2) $\dfrac{\sqrt2}{12}$
  \item[\textbf{20.}] (1) $a_n=\begin{cases}120,&n=1,\\130-10n,&2\leqslant n\leqslant6,\\70\left(\dfrac34\right)^{n-6},&n\geqslant7.\end{cases}$；(2) 证明见解析
  \item[\textbf{21.}] (1) $y^2=4x$（$x\geqslant0$）或 $y=0$（$x<0$）；(2) 最小值 $0$
  \item[\textbf{22.}] (1) $a\leqslant2$ 时在 $(0,+\infty)$ 上递增；$a>2$ 时在 $(0,x_1)$、$(x_2,+\infty)$ 上递增，在 $(x_1,x_2)$ 上递减，其中 $x_{1,2}=\dfrac{a\mp\sqrt{a^2-4}}2$；(2) 不存在
\end{itemize}
